%% toc: P2 9 9.5,9.6,9.7
\chapter{Aristotle's Fourteen Syllogism Rules under the Empty-Term Policy}
\label{sec:p02-w02-p02-002-fourteen}

The calculi \(\mathsf S\) and \(\mathsf S^\dagger\) of
Chapter~\ref{sec:p02-syllogistic-inference} let a predicate extension be
empty. George Boger's reconstruction of Aristotle's underlying logic, presented in
Chapter~\ref{sec:p13-src063-0103}, assigns existential import to universal
affirmations. This chapter tests that difference on one fixed inventory: the
fourteen two-premise syllogism rules in Boger's summary
\cite[p.~188, table~20]{Boger2004UnderlyingLogic}. It states which rules hold
in \(\mathsf S\) without an existence premise and which do not. For each
failure it gives a countermodel with an empty term and the weakest existence
premise that restores the rule. It also locates the one step of the
historical completion proofs at which existence enters.
Chapter~\ref{sec:p02-w02-p02-001-historical-contrast} compares the modern
calculi with the two historical reconstructions assumption by assumption.
The results below specialize that comparison to the fourteen rules and do not
repeat it.

\emph{Inputs.} The language, set semantics and rules of \(\mathsf S\)
(Definitions~\ref{def:p02-syllogistic-language},
\ref{def:p02-syllogistic-semantics}
and~\ref{def:p02-syllogistic-identifications}
and the rule sheet of Section~\ref{sec:p02-syllogistic-inference}); soundness and completeness of \(\mathsf S\)
(Theorem~\ref{thm:p02-syllogistic-s-completeness}); the \(\mathcal S\)-reading
\(\sigma\) and the existence-committing evaluation \(\models_{\mathrm B}\) of
Definition~\ref{def:p02-w02-p02-001-patterns}, built from the conditions
\(\alpha,\iota,\varepsilon,\omega\) of
Definition~\ref{def:p13-boger-guarded-bridge}; and the rule inventory
SR1--SR14 of Section~\ref{sec:p13-boger-deduction-rules}.

\emph{Conclusions.} Under the empty-term policy of \(\mathsf S\), twelve of
the fourteen rules are derivable by one rule application each, possibly
after a formula identification (Theorem~\ref{thm:p02-w02-p02-002-twelve}).
Darapti and Felapton fail. Each holds exactly when the premise that the
middle term is nonempty is added, and the resulting three-premise argument
has no superfluous premise (Theorem~\ref{thm:p02-w02-p02-002-third-figure}).
Under the existence-committing evaluation all fourteen rules are valid
(Theorem~\ref{thm:p02-w02-p02-002-committing}). Of the three conversions,
only conversion per accidens depends on existence
(Lemma~\ref{thm:p02-w02-p02-002-conversions}), and it is the one step at
which Boger's completions of Darapti and Felapton use existence
(Method~\ref{def:p02-w02-p02-002-locate}).

\section{The fourteen rules and their reading in \texorpdfstring{\(\mathcal S\)}{S}}
\label{sec:p02-w02-p02-002-reading}

Boger reads the \emph{Prior Analytics} as identifying fourteen panvalid
two-premise patterns in three figures. He contrasts this with the
twenty-four syllogisms in four figures counted by later traditional logic
\cite[p.~184 and note~75]{Boger2004UnderlyingLogic}. In his reconstruction
the deduction system consists of four kinds of categorical sentence, the
contradictory and contrary pairs, three conversion rules, the fourteen
syllogism rules, and direct and indirect deduction
\cite[p.~166, table~10]{Boger2004UnderlyingLogic}. Aristotle limited the
study of multi-premise arguments to their two-premise patterns
\cite[p.~160]{Boger2004UnderlyingLogic}. The fourteen rules and the three
conversions are displayed in Section~\ref{sec:p13-boger-deduction-rules}.

\begin{hbconvention}{def:p02-w02-p02-002-reading}{Term letters as atoms}
A categorical sentence in Boger's notation has the form \(XyZ\), where
\(y\) is one of the constants \(a,e,i,o\) and \(X\), the predicate, and
\(Z\), the subject, are distinct term letters. For each term letter \(T\)
among \(A,B,C,M,N,X,P,R,S\), let the lower-case letter \(t\) name an atom,
with distinct letters naming distinct atoms. Thus \(m\) and \(n\) below are
atoms, not literal metavariables. The sentence \(XyZ\) is then the pattern
\(x\mathbin{\mathsf y}z\) of Definition~\ref{def:p02-w02-p02-001-patterns},
and it has the \(\mathcal S\)-reading
\[
 \sigma(XaZ)=\forall(z,x),\qquad
 \sigma(XeZ)=\forall(z,\bar x),
\]
\[
 \sigma(XiZ)=\exists(z,x),\qquad
 \sigma(XoZ)=\exists(z,\bar x)
\]
and the existence-committing truth condition \(\models_{\mathrm B}\) of the
same definition. A rule \(\varphi_1,\varphi_2\mathbin{/}\psi\) \emph{holds in
\(\mathsf S\)} when
\(\{\sigma(\varphi_1),\sigma(\varphi_2)\}\vdash_{\mathsf S}\sigma(\psi)\), and
\emph{holds under \(\models_{\mathrm B}\)} when
\(\{\varphi_1,\varphi_2\}\models_{\mathrm B}\psi\).
\end{hbconvention}

The subject comes second in \(XyZ\) and first in \(\sigma(XyZ)\). Every
\(\sigma(XyZ)\) is an \(\mathcal S\)-formula, because its first literal is an
atom. In each of the fourteen rules the three term letters are distinct, so
the three atoms are distinct.

\begin{longtable}{@{}>{\raggedright\arraybackslash}p{0.19\linewidth}>{\raggedright\arraybackslash}p{0.22\linewidth}>{\raggedright\arraybackslash}p{0.15\linewidth}>{\raggedright\arraybackslash}p{0.34\linewidth}@{}}
\caption{The fourteen rules read in \(\mathcal S\)}
\label{tab:p02-w02-p02-002-fourteen}\\
\hline
Rule & Premises read & Conclusion read & Derivation in \(\mathsf S\) \\
\hline
\endfirsthead
\hline
Rule & Premises read & Conclusion read & Derivation in \(\mathsf S\) \\
\hline
\endhead
SR1 Barbara & \(\forall(b,a)\), \(\forall(c,b)\) & \(\forall(c,a)\) & B. \\
SR2 Celarent & \(\forall(b,\bar a)\), \(\forall(c,b)\) & \(\forall(c,\bar a)\) & B. \\
SR3 Darii & \(\forall(b,a)\), \(\exists(c,b)\) & \(\exists(c,a)\) & D1. \\
SR4 Ferio & \(\forall(b,\bar a)\), \(\exists(c,b)\) & \(\exists(c,\bar a)\) & D1. \\
SR5 Cesare & \(\forall(n,\bar m)\), \(\forall(x,m)\) & \(\forall(x,\bar n)\) & Read \(\forall(n,\bar m)\) as \(\forall(m,\bar n)\); B. \\
SR6 Camestres & \(\forall(n,m)\), \(\forall(x,\bar m)\) & \(\forall(x,\bar n)\) & Read \(\forall(x,\bar m)\) as \(\forall(m,\bar x)\); B gives \(\forall(n,\bar x)\), read as \(\forall(x,\bar n)\). \\
SR7 Festino & \(\forall(n,\bar m)\), \(\exists(x,m)\) & \(\exists(x,\bar n)\) & Read \(\forall(n,\bar m)\) as \(\forall(m,\bar n)\); D1. \\
SR8 Baroco & \(\forall(n,m)\), \(\exists(x,\bar m)\) & \(\exists(x,\bar n)\) & D3. \\
SR9 Darapti & \(\forall(s,p)\), \(\forall(s,r)\) & \(\exists(r,p)\) & None: fails when \(s\) is empty. \\
SR10 Felapton & \(\forall(s,\bar p)\), \(\forall(s,r)\) & \(\exists(r,\bar p)\) & None: fails when \(s\) is empty. \\
SR11 Disamis & \(\exists(s,p)\), \(\forall(s,r)\) & \(\exists(r,p)\) & D2. \\
SR12 Datisi & \(\forall(s,p)\), \(\exists(s,r)\) & \(\exists(r,p)\) & Read \(\exists(s,r)\) as \(\exists(r,s)\); D1. \\
SR13 Bocardo & \(\exists(s,\bar p)\), \(\forall(s,r)\) & \(\exists(r,\bar p)\) & D2. \\
SR14 Ferison & \(\forall(s,\bar p)\), \(\exists(s,r)\) & \(\exists(r,\bar p)\) & Read \(\exists(s,r)\) as \(\exists(r,s)\); D1. \\
\hline
\end{longtable}

\begin{hblemma}{thm:p02-w02-p02-002-conversions}{proved=annex:p02-w02-p02-002-conversions}{The three conversions in \(\mathsf S\)}
Let \(A,B\) be distinct term letters.
\begin{enumerate}
\item \(\sigma(AeB)=\forall(b,\bar a)\) and \(\sigma(BeA)=\forall(a,\bar b)\)
 are identified by universal contraposition, so each is derivable from the
 other in \(\mathsf S\) with no rule application.
\item \(\sigma(AiB)=\exists(b,a)\) and \(\sigma(BiA)=\exists(a,b)\) are
 identified by existential symmetry, so each is derivable from the other
 with no rule application.
\item \(\forall(b,a)\nvdash_{\mathsf S}\exists(a,b)\): conversion per
 accidens, from \(AaB\) to \(BiA\), does not hold in \(\mathsf S\). It holds
 once \(\exists(b,b)\) is added:
 \(\{\forall(b,a),\exists(b,b)\}\vdash_{\mathsf S}\exists(a,b)\).
\end{enumerate}
\end{hblemma}

Boger's conversion rules are stated in his table~16
\cite[p.~183]{Boger2004UnderlyingLogic}; their justification under the
existent-subject policy is Lemma~\ref{thm:p13-boger-conversions}. Item~3
shows that in \(\mathsf S\) the existence of the subject of the universal
premise must be supplied as a separate premise.

\begin{hbtheorem}{thm:p02-w02-p02-002-twelve}{proved=annex:p02-w02-p02-002-twelve}{Twelve rules need no existence premise}
For every rule SR\(k\) with \(k\notin\{9,10\}\), the premises' readings
derive the conclusion's reading in \(\mathsf S\) by the single rule
application listed in Table~\ref{tab:p02-w02-p02-002-fourteen}, together
with the formula identifications listed there. Consequently
\(\{\sigma(\varphi_1),\sigma(\varphi_2)\}\models\sigma(\psi)\) for each of
these twelve rules \(\varphi_1,\varphi_2\mathbin{/}\psi\), with no
restriction on which predicate extensions are empty.
\end{hbtheorem}

\section{Countermodels in the third figure}
\label{sec:p02-w02-p02-002-countermodels}

\begin{hbtheorem}{thm:p02-w02-p02-002-third-figure}{proved=annex:p02-w02-p02-002-third-figure}{Exactly which existence premise Darapti and Felapton need}
Put
\[
 \Gamma_9=\{\forall(s,p),\forall(s,r)\},\qquad \theta_9=\exists(r,p),
\]
\[
 \Gamma_{10}=\{\forall(s,\bar p),\forall(s,r)\},\qquad
 \theta_{10}=\exists(r,\bar p),
\]
the readings of Darapti and Felapton. Let \(k\in\{9,10\}\) and
\[
 E\subseteq\{\exists(p,p),\exists(r,r),\exists(s,s)\}.
\]
\begin{enumerate}
\item \(\Gamma_k\cup E\vdash_{\mathsf S}\theta_k\) if and only if
 \(\exists(s,s)\in E\). In particular
 \(\Gamma_k\nvdash_{\mathsf S}\theta_k\), and the premises that the major
 and minor terms are nonempty do not restore the rule.
\item No proper subset of \(\Gamma_k\cup\{\exists(s,s)\}\) has
 \(\theta_k\) as a semantic consequence.
\item \(\theta_k\) is neither a member of \(\Gamma_k\cup\{\exists(s,s)\}\)
 nor identified with a member of it by
 Definition~\ref{def:p02-syllogistic-identifications}.
\end{enumerate}
\end{hbtheorem}

The term whose existence is needed is the middle term \(S\). It is the
subject of both premises and does not occur in the conclusion. Supplying a
witness for the conclusion's own terms \(R\) and \(P\) does not help, as
item~1 shows.

\begin{hbtheorem}{thm:p02-w02-p02-002-committing}{proved=annex:p02-w02-p02-002-committing}{The fourteen rules under existential import}
For every \(k\in\{1,\ldots,14\}\), the rule SR\(k\) holds under
\(\models_{\mathrm B}\). For SR9 and SR10, every structure satisfying the
premises under \(\models_{\mathrm B}\) satisfies \(\exists(s,s)\).
Consequently the empty-term policy of \(\mathsf S\) and the
existence-committing evaluation disagree on exactly two of the fourteen
rules, Darapti and Felapton.
\end{hbtheorem}

\begin{hbmethod}{def:p02-w02-p02-002-locate}{Locating the existence step in a completion}
\emph{Input.} One of Boger's completion deductions for a second- or
third-figure rule, in his table~19
\cite[p.~187]{Boger2004UnderlyingLogic}, also given in
Theorem~\ref{thm:p13-boger-ten-completions}.
\emph{Output.} A direct \(\mathsf S\)-derivation of the rule's reading, or
the exact premise that such a derivation needs.
\begin{enumerate}
\item Replace each line by its reading under \(\sigma\).
\item Delete every e-conversion and i-conversion step. By
 Lemma~\ref{thm:p02-w02-p02-002-conversions}, items 1 and~2, its premise and
 conclusion are identified formulas.
\item Replace every application of Barbara or Celarent by an instance of
 B, and every application of Darii or Ferio by an instance of D1, with the
 term letters of that application in place of those in rows SR1--SR4 of
 Table~\ref{tab:p02-w02-p02-002-fourteen}.
\item Replace an indirect deduction, as in Baroco and Bocardo, by the
 direct D3 or D2 derivation of Table~\ref{tab:p02-w02-p02-002-fourteen}.
 \(\mathsf S\) has no rule that discharges an assumption.
\item For an a-conversion step from \(RaS\) to \(SiR\), add the premise
 \(\exists(s,s)\) and replace the step by D1 applied to \(\forall(s,r)\)
 and \(\exists(s,s)\). By Lemma~\ref{thm:p02-w02-p02-002-conversions},
 item~3, without that premise the step has no \(\mathsf S\)-counterpart.
\end{enumerate}
\emph{Side conditions.} The three term letters are distinct atoms. The
universe is nonempty, as in every structure of
Definition~\ref{def:p02-syllogistic-semantics}; the method adds no other
existence assumption. In table~19 the a-conversion step occurs only in the
completions of Darapti and Felapton, as line~4, from \(RaS\) to \(SiR\).
Step~5 is therefore needed for these two rules and for no other.
\end{hbmethod}

\begin{hbexample}{def:p02-w02-p02-002-felapton}{Felapton with an empty middle term}
Let \(s\), \(r\) and \(p\) mean \emph{student scoring above 100 percent},
\emph{student} and \emph{prize winner}. The premises \(PeS\) and \(RaS\)
say that no student scoring above 100 percent wins a prize and that every
such student is a student. Felapton's conclusion \(PoR\) says that some
student does not win a prize.

\emph{Countermodel in \(\mathsf S\).} Take the universe \(\{k\}\), with
\(s^{\mathfrak A}=\varnothing\) and
\(r^{\mathfrak A}=p^{\mathfrak A}=\{k\}\): one student, who wins a prize,
and nobody above 100 percent. Then \(\forall(s,\bar p)\) and
\(\forall(s,r)\) are true because their subject is empty, and
\(\exists(r,\bar p)\) is false because \(r^{\mathfrak A}\setminus
p^{\mathfrak A}=\varnothing\). The same structure satisfies \(\exists(r,r)\)
and \(\exists(p,p)\), so the existence of students and of prize winners does
not rescue the inference.

\emph{Restored derivation.} Add \(\exists(s,s)\). D1 applied to
\(\forall(s,r)\) and \(\exists(s,s)\) gives \(\exists(s,r)\), the
counterpart of Boger's line \(SiR\). D1 applied to \(\forall(s,\bar p)\) and
\(\exists(s,s)\) gives \(\exists(s,\bar p)\), and D2 applied to
\(\exists(s,\bar p)\) and \(\forall(s,r)\) gives \(\exists(r,\bar p)\).

\emph{Existence-committing reading.} In the countermodel \(RaS\) is false
under \(\models_{\mathrm B}\), since \(\alpha(\varnothing,\{k\})\) fails. The
structure is therefore no counterexample to Felapton under that reading, in
agreement with Theorem~\ref{thm:p02-w02-p02-002-committing}.
\end{hbexample}

\begin{hbwarning}{def:p02-w02-p02-002-limits}{What these verdicts do and do not compare}
The verdicts compare two extensional evaluations on fourteen fixed
patterns. They do not test Boger's intensional or possessional reading of
predication, which Definition~\ref{def:p13-boger-guarded-bridge} and
Warning~\ref{def:p02-w02-p02-001-limits} exclude from the comparison.
Theorem~\ref{thm:p02-w02-p02-002-third-figure} restores Darapti and Felapton
only as three-premise arguments in \(\mathcal S\). A three-premise argument
is not one of the two-premise patterns to which the reconstruction confines
the syllogism rules \cite[pp.~160, 166]{Boger2004UnderlyingLogic}. A rule
that holds under \(\models_{\mathrm B}\) is not thereby a rule of
\(\mathsf S\): an inference valid under \(\models_{\mathrm B}\) may need an
existence premise, or a case split over partial privatives, before
\(\mathsf S\) can derive it
(Theorem~\ref{thm:p02-w02-p02-001-reduction}).
\end{hbwarning}

\section{Relation to the two historical reconstructions}
\label{sec:p02-w02-p02-002-historical}

The calculi \(\mathsf S\) and \(\mathsf S^\dagger\) are not Aristotle's
system, and neither is identified here with either historical
reconstruction. The assumptions in which the results of this chapter
differ from Boger's reconstruction, presented in
Chapter~\ref{sec:p13-src063-0103}, are as follows.
\begin{itemize}
\item \emph{Existence of subjects.} In Boger's semantics a sentence about an
 existent subject is true or false as the state of affairs obtains, and in
 these cases there are no empty classes; for a nonexistent subject, every
 affirmation is false and every sentence with a negative operator true
 \cite[p.~159]{Boger2004UnderlyingLogic}. The ordinary derivations of
 Section~\ref{sec:p13-boger-deduction} assume an existent subject matter
 (Definition~\ref{def:p13-boger-deduction-domain}). In \(\mathsf S\) a
 predicate extension may be empty, and Darapti and Felapton fail
 (Theorem~\ref{thm:p02-w02-p02-002-third-figure}).
\item \emph{Conversion per accidens.} Boger's table~16 states
 a-conversion as a rule \cite[p.~183]{Boger2004UnderlyingLogic}. In
 \(\mathsf S\) it requires the extra premise \(\exists(b,b)\)
 (Lemma~\ref{thm:p02-w02-p02-002-conversions}).
\item \emph{Indirect deduction.} Boger's completions of Baroco and Bocardo
 assume the contradictory of the conclusion and discharge it
 \cite[p.~187, table~19]{Boger2004UnderlyingLogic}. The \(\mathsf S\)
 derivations of these rules are direct, by D3 and D2.
\item \emph{Two premises.} The syllogism rules are two-premise patterns
 \cite[pp.~160, 166]{Boger2004UnderlyingLogic}. In \(\mathsf S\) the
 counterparts of Darapti and Felapton need three premises.
\item \emph{Inventory.} The reconstruction counts fourteen rules in three
 figures \cite[p.~184]{Boger2004UnderlyingLogic}. \(\mathsf S\) has no
 figures and no list of moods. By completeness, every semantically valid
 inference between \(\mathcal S\)-formulas is derivable, including the
 twelve rules of Theorem~\ref{thm:p02-w02-p02-002-twelve}.
\end{itemize}

John Woods and Andrew Irvine's reconstruction of Aristotle's early logic
\cite[pp.~27--99]{WoodsIrvine2004EarlyLogic} has no presentation in this
handbook and is cited directly. Its differences from \(\mathsf S\) and
\(\mathsf S^\dagger\) are stated in
Section~\ref{sec:p02-w02-p02-001-woods-irvine}. Two of its conditions bear
on the three-premise arguments of
Theorem~\ref{thm:p02-w02-p02-002-third-figure}. Under their condition Min,
a syllogism contains the premises its validity needs and no others; under
Non-Circ, its conclusion repeats no premise
\cite[p.~54]{WoodsIrvine2004EarlyLogic}. Items 2 and~3 of
Theorem~\ref{thm:p02-w02-p02-002-third-figure} are formal counterparts of
these two conditions for the arguments
\(\Gamma_k\cup\{\exists(s,s)\}\mathbin{/}\theta_k\), with semantic
consequence in place of necessitation and the identifications of
Definition~\ref{def:p02-syllogistic-identifications} in place of
synonymy. Their third condition, Prem+, is a condition Woods and Irvine describe
as expressly proclaimed in the \emph{Prior Analytics}
\cite[p.~63]{WoodsIrvine2004EarlyLogic}; its exact content is given in
Section~\ref{sec:p02-w02-p02-001-woods-irvine}. These counterparts do not
show that Woods and Irvine would classify the three-premise arguments as
syllogisms. Their notion of validity and its relation to existence are
compared with \(\mathsf S\) in that section, not here.
\index{syllogistic logic!existential import}
\index{syllogistic logic!fourteen syllogism rules}
\index{existential import!Darapti and Felapton}
\index{conversion per accidens}
